\chapter{Objetivos}
\label{ch:objetivos}

\section{Objetivo geral}
\label{sec:objetivo-geral}

Desenvolver, implementar e documentar o \CabotageLens{} como um framework computacional auditável para comparar o transporte rodoviário direto com uma cadeia rodoviária--cabotagem--rodoviária em corredores brasileiros, considerando emissões operacionais \ttw{} de \cooe{} e proxy de custo operacional modelado sob fronteiras metodológicas explícitas.

\section{Objetivos específicos}
\label{sec:objetivos-especificos}

Os objetivos específicos do trabalho são:

\begin{itemize}
  \item Definir a unidade funcional do estudo e as alternativas comparáveis como cadeias porta a porta, distinguindo a rota rodoviária direta da cadeia rodoviária--cabotagem--rodoviária.
  \item Construir alternativas orientadas por rota para pares origem-destino selecionados, incluindo acesso rodoviário inicial, perna marítima, operações portuárias e \emph{hoteling} quando aplicáveis, e acesso rodoviário final.
  \item Reconstruir as viagens de cabotagem observadas nas tabelas de Carga e Atracação da ANTAQ, preservando a ordem das escalas, os embarques e desembarques e a carga a bordo em cada subtrecho.
  \item Considerar, para cada par portuário direcionado, todas as viagens em que a origem antecede o destino, incluindo ligações diretas, escalas intermediárias e corredores distintos observados na base.
  \item Associar as viagens aos indicadores de intensidade por trabalho de transporte do EU MRV por meio do número IMO e resolver ausências com estatística robusta de classe ou tipo de navio, registrando a origem de cada valor.
  \item Estimar a intensidade marítima representativa de cada par origem--destino pela mediana ponderada pelo trabalho de transporte das viagens elegíveis, mantendo a escolha do corredor físico separada da agregação estatística da intensidade.
  \item Preservar a proveniência de distâncias, seleção de portos, construção da perna marítima e insumos de operações portuárias/\emph{hoteling}, incluindo níveis de fonte, \emph{fallbacks}, avisos e dados indisponíveis.
  \item Estimar, para os cenários selecionados, emissões operacionais \ttw{} de \cooe{} e proxy de custo operacional modelado, mantendo unidades, fronteiras e limites de interpretação explícitos.
  \item Expor hipóteses, aproximações, avisos de rota, lacunas de dados e componentes indisponíveis, evitando que ausência de informação seja interpretada como valor nulo.
  \item Documentar a implementação computacional com ênfase em rastreabilidade, reprodutibilidade e recursos acessíveis, priorizando bibliotecas de código aberto e serviços públicos, gratuitos ou com camada gratuita quando aplicável, de modo a facilitar auditoria, inspeção e extensão futura do protótipo.
  \item Classificar os resultados de forma conservadora segundo qualidade da evidência, sensibilidade a premissas, limitações de rota e comparabilidade externa.
  \item Usar análises de sensibilidade e evidência externa de benchmark para apoiar interpretação acadêmica defensável, sem afirmar equivalência calibrada de magnitude.
  \item Discutir limitações e melhorias futuras para o uso metodológico, acadêmico e computacional do \CabotageLens{}.
\end{itemize}

Esses objetivos delimitam a leitura do relatório. O trabalho promete uma comparação auditável e condicionada por fronteiras explícitas; não promete resolver contratação comercial, disponibilidade operacional, preço de mercado, super-rede multimodal completa ou superioridade universal da cabotagem.
